\section{Primes in geometric intervals}
Barthel and M\"uller ask whether every interval $[q^t,q^{t+1})$ meets every
reduced residue class modulo $q$, for integers $q\ge2$ and $t\ge1$.
Their Lemmas 5--6 cover all $q\le45000$ (with Bertrand's postulate for $q=2$).
Their Remark 7 obtains all $t\ge3$ on Dirichlet GRH, and $t=2$ for
$q\ge17386763$. The following closes the entire intervening range. The
finite computation also extends their $t=1$ verification.
\hfill\href{https://doi.org/10.1080/00029890.2022.2116250}{[Barthel--M\"uller, 2022]}

\paragraph{Theorem.} For every $q\ge2$, every $a$ coprime to $q$, and
every integer $t\ge2$, some prime $p\equiv a\pmod q$ satisfies
$q^t<p<q^{t+1}$. The same conclusion holds for $t=1$ whenever $q\le50000$.

\paragraph{Proof.}
Dudek, Greni\'e and Molteni, Theorem 1.1 and Table 1, prove under Dirichlet
GRH that there is such a prime within distance
\[
 H(x)=q\sqrt x\bigl(0.6265\log x+\log q+14\bigr)
\]
of $x$ if $q\ge3$ and $x\ge(18q\log q)^2$. We have enlarged both their
radius and their threshold using $\varphi(q)\le q$. Their explicit formula
applies to primitive Dirichlet functions and restores the factors at primes
dividing the modulus. This is precisely the zero hypothesis exported by the
Lean repository.\hfill\href{https://doi.org/10.1142/S1793042119500441}{[Dudek--Greni\'e--Molteni, 2019]}

For $q\ge50000$ write $L=\log q$ and $x_2=(18qL)^2$.
The quotient $q/(378L^2+252L)$ is increasing for $L>2$ by
differentiation. At $q=50000$, $10<L<11$, giving
$378L^2+252L<48510<50000$. Thus $x_2<q^3$ and
\[
 \frac{H(x_2)}{q^2}<18L(3L+14),\qquad
 q^2<x_2-H(x_2)<x_2+H(x_2)<q^3.
\]
Indeed the upper endpoint divided by $q^2$ is at most $378L^2+252L$,
and the lower one exceeds $270L^2-252L>1$. For $t\ge2$ set
$x_t=q^{t-2}x_2$. The function $H(x)/x$ decreases for $x\ge x_2$,
again by differentiation. Therefore the interval of radius $H(x_t)$
about $x_t$ is contained in the dilation by $q^{t-2}$ of the interval
just proved to lie inside $(q^2,q^3)$. This proves every $t\ge2$ in the
large-modulus range.

For $45001\le q\le50000$ the included exact sieve verifies $t=1,2$
in every reduced class. The following numbers are maxima of the
\emph{first} prime witnesses when the search starts at the indicated
left endpoint. The endpoints were chosen to let every modulus share a
single sieve.
\begin{center}
\begin{tabular}{rrrr}
$t$ & Search starts at & Classes checked & Largest witness\\ \hline
1 & $50001$ & $144390310$ & $9819037$\\
2 & $2500000001$ & $144390310$ & $2515292327$
\end{tabular}
\end{center}
For each row the program sieves $25000000$ consecutive integers by all
primes through the square root of the last integer. It verifies that the
segment lies inside the appropriate geometric interval for each $q$,
then traverses each reduced progression until it finds a prime. This
uses only integer arithmetic. For these moduli $t\ge3$ follows from
Barthel--M\"uller's Remark 7. Their small-modulus results complete the proof.
Neither endpoint can be a prime in the indicated class for $q\ge2$.

\section{An exact obstruction to Pomerance's divisibility conjecture}
For a prime $p\ge3$, put
\[
 B(p)=\prod_{r<p,\ r\ {\rm prime}}r(r-1),\qquad
 D(p)=\frac{p-1}{\gcd(p-1,B(p))}.
\]
Pomerance conjectured $D(p)=1$ for every $p$. He observed that if the
least prime $P(q)\equiv1\pmod q$ satisfies $q^2\mid P(q)-1$ for a
prime $q$, then $p=P(q)$ violates the conjecture. The converse also
holds when we use the repository's Dirichlet GRH. The proof gives the
exact size of every possible divisibility defect.
\hfill\href{https://math.dartmouth.edu/~carlp/PDF/carmichaelconjecture.pdf}{[Pomerance, 1974]}

\paragraph{Theorem.} Either $D(p)=1$, or $D(p)$ equals a single
prime $q$ for which $p=P(q)$ and $v_q(p-1)=2$. In particular the
conjecture is equivalent to the assertion that
$q^2\nmid P(q)-1$ for every prime $q$.

\paragraph{Proof.}
Fix a prime $q\mid p-1$ and put $e=v_q(p-1)$. The factor $q$ itself
occurs in $B(p)$. If $q\le50000$, the geometric-interval theorem for
$t=1,\ldots,e-1$ supplies $e-1$ distinct primes $r<p$ with
$q\mid r-1$. This handles all $e$, also when the list is empty.

Suppose $q>50000$. Above the class $1\pmod q$ there are two reduced
classes modulo $3q$, by the Chinese remainder theorem. The explicit
least-prime theorem of Lamzouri, Li and Soundararajan (Corollary 1.2)
gives a prime in each such class at most
$[2(q-1)\log(3q)]^2<q^3$. The final inequality follows from
$4\log^2(3q)<q$, valid at $50000$ and thereafter by monotonicity of
$q/\log^2(3q)$. These primes are distinct. When $e\ge3$, use these two
primes first and then one prime from each of the $e-3$ intervals
$[q^j,q^{j+1})$, $j=3,\ldots,e-1$, always in class $1\pmod q$.
They all lie below $q^e\le p-1$. Together with the factor $q$ they give
$v_q(B(p))\ge e$.\hfill\href{https://arxiv.org/pdf/1309.3595}{[Lamzouri--Li--Soundararajan, Cor.~1.2]}

The only possible shortage is $e=2$. It occurs precisely when there
is no earlier prime $r\equiv1\pmod q$, i.e., when $p=P(q)$, and the
shortage is then exactly one. If $q<s$ were two deficient primes at
$p$, we would have
$q^2s^2\le p-1<p=P(q)\le[(q-1)\log q]^2$ by the same
least-prime theorem. This forces $s<\log q<q$, a contradiction.

\paragraph{Frequency and a finite exclusion.}
Every exceptional $p\le X$ has a distinct associated $q$, because $p$
is determined by $q$. Both $p$ and $q$ are odd and $q^2\mid p-1$,
so $2q^2\le p-1$. Thus there are at most
$\pi(\sqrt{(X-1)/2})$ exceptions, in particular
$O(\sqrt X/\log X)$. The relative density of exceptions among primes
is zero. An independent finite check in the archive finds $P(q)<q^2$
for all $664579$ prime moduli $q\le10^7$. It sieves the moduli and
tests candidates $kq+1$ in increasing order by exact trial division.
The largest least prime encountered is $1753933943$, and the largest
multiplier $k$ is $234$, first attained at $q=4722079$. In fact the finite bound
\[
 p\le 2\cdot10^{14}+1\qquad\Longrightarrow\qquad D(p)=1
\]
can be certified independently of GRH. For every prime $q\le10^7$ and every
$e\ge2$ with $q^e\le2\cdot10^{14}$, the same program verifies that the
$(e-1)$st prime in $1\bmod q$ lies below $q^e$. This checks $671350$
prime powers. If $q$ is odd and $q^2\mid p-1$, then $2q^2\le p-1$,
so these ranges include all relevant $q,e$ for the asserted bound, also
for $q=2$. The factor $q$ and the $e-1$ earlier primes supply $q^e$
in $B(p)$. A counterexample would require $q>10^7$ and
$p=1+2kq^2\le[(q-1)\log q]^2$ for some positive integer $k$,
with every smaller candidate prime in $1\pmod q$ absent.

The universally quantified $t=1$ question, the existence of such a $q$,
and Carmichael's conjecture about multiplicities of totients all remain
unresolved by these arguments. Pomerance's divisibility assertion would
rule out one specific proposed construction of a totient with a single
preimage. It would not by itself settle the multiplicity conjecture.
